\chapter*{مقدمة}
\addcontentsline{toc}{schapter}{المقدمة}

في عام 1992 اقترح هي (He) طريقة التغاير التكراري (VIM) [1]. وتستخدم هذه الطريقة حالياً على نطاق واسع من قبل العديد من الباجثين لدراسة المعادلات التفاضلية الجزئية الخطية وغير الخطية. وتعد الطريقة عملية موثوقة وفعالة لمجموعة واسعة من التطبيقات العلمية والهندسية ، سواء كانت خطية او غير خطية ، متجانسة او غير متجانسة ، او لأنظمة المعادلات ايضاً ، وقد اظهرت قدرتها من خلال العديد من المؤلفين [2] انها طريقة اكثر قوة من التقنيات الأخرى مثل طريقة تحليل ادوميان [3] وطريقة الاضطراب الهوموتبي وغيرها.\\

توفر هذه الطريقة تقارباً سريعاً للتقريبات المتتالية للحل الدقيق عندما يكون الحل موجوداً؛ وإلا فيمكن استخدام عدد قليل من التقريبات للأغراض العددية ، وتختلف التقنيات العددية الحالية عن طريقة VIM بسبب الافتراضات التقييدية المستخدمة للتعامل مع الحدود غير الخطية ، اما طريقة VIM فلا تتطلب شروطاً خاصة مثل الخطية او المعاملات الصغيرة او كثيرات حدود ادوميان للتعامل مع المؤثرات غير الخطية ، ومن المزايا المهمة الاخرى ان طريقة VIM قادرة على تقليل حجم الحسابات بشكل كبير مع الحفاظ على دقة عالية للحل العددي، اضافة الى ذلك فأن قوة الطريقة تمنحها قابلية للتطبيق في معالجة عدد كبير من المسائل التحليلية والعددية. [5]\\

تم انجاز قدر كبير من الابحاث في دراسة انظمة المعادلات التفاضلية الجزئية الخطية وغير الخطية (PDEs) ، وتنشأ انظمة المعادلات التفاضلية الجزئية الخطية وغير الخطية في العديد من النماذج العلمية ، مثل انتشار امواج المياه الضحلة ونموذج "بروسليتور" (Brusselator) في تفاعلات الانتشار الكيميائي.[7]